\par\medskip\noindent
2.6. En 2.5 solo se hizo variar $G$, manteniendo fijos $P$ y $Q$.
Conviene ahora decir algunas palabras acerca de la variación en $P,Q$.
El hecho esencial, del cual los demás se deducen formalmente, es que la
dependencia contravariante de $\operatorname{BIEXT}(P,Q;G)$ respecto
de $P,Q$ \underline{conmuta} con su dependencia covariante respecto de
$G$. Por tanto, las categorías consideradas pueden verse

%%%%%%%%%%  scan idx 180  |  folio 169  %%%%%%%%%%

como las categorías fibra de una \underline{categoría cofibrada}
\begin{equation}\tag{2.6.1}
 \operatorname{BIEXT}'(\underline T)
 \longrightarrow
 \operatorname{GRAB}(\underline T)^{\mathrm{op}}\times
 \operatorname{GRAB}(\underline T)^{\mathrm{op}}\times
 \operatorname{GRAB}(\underline T),
\end{equation}
donde los superíndices $\mathrm{op}$ designan las categorías opuestas y
el funtor estructural es $E\mapsto(P,Q,G)$. Esta categoría cofibrada es
aditiva en cada uno de sus tres argumentos o, como diremos,
\underline{triaditiva}.

Parafraseando argumentos estándar, que no reproduciremos, se deduce
que la composición $E.E'$ en la categoría fibra
\sisi{$\operatorname{BIEXT}(P,P';G)$}%
      {$\operatorname{BIEXT}(P,Q;G)$}
puede describirse usando la estructura aditiva de cualquiera de los
tres argumentos: las tres estructuras de pseudogrupo resultantes son
canónicamente isomorfas. En consecuencia, los grupos abelianos
\begin{equation}\tag{2.6.2}
 \operatorname{Biext}^{i}(P,Q;G),\qquad i=0,1,
\end{equation}
pueden considerarse como dos funtores trilineales en $P,Q,G$,
respectivamente en $\operatorname{GRAB}(\underline T)^{\mathrm{op}}$,
$\operatorname{GRAB}(\underline T)^{\mathrm{op}}$ y
$\operatorname{GRAB}(\underline T)$. Equivalentemente, son
\underline{funtores trilineales} de los argumentos variables $P,Q,G$
en $\operatorname{GRAB}(\underline T)$, contravariantes en $P,Q$ y
covariantes en $G$, con valores en la categoría de grupos abelianos.
Con esta varianza, el isomorfismo (2.5.4.1) es funtorial en todos los
argumentos $P,Q,G$, y no solo en $G$. La afirmación análoga será válida
más adelante para la descripción de $\operatorname{Biext}^{1}$
mediante un $\operatorname{Ext}^{1}$.

2.7. \underline{Biextensión simétrica de una biextensión dada}.

Si $E$ es una biextensión de $(P,Q)$ por $G$, definamos su
\underline{biextensión simétrica} $\,{}^{s}E$, que es una biextensión
de $(Q,P)$ por $G$. El haz de conjuntos subyacente a $\,{}^{s}E$ es el
de $E$, y las acciones de $G$ sobre $E$ y $\,{}^{s}E$ son las mismas.
El morfismo estructural
% Canon: SGA-SRC-001256.
\[
 {}^{s}E\longrightarrow \sisi{P\times Q}{Q\times P}
\]
es el compuesto de $E\longrightarrow P\times Q$ con la simetría
$P\times Q\longrightarrow Q\times P$; finalmente, las estructuras de

%%%%%%%%%%  scan idx 181  |  folio 170  %%%%%%%%%%

extensión izquierda y derecha de $\,{}^{s}E$ son, respectivamente,
las estructuras de extensión derecha e izquierda de $E$. Es evidente
que $E\mapsto{}^{s}E$ es un automorfismo de la categoría
$\operatorname{BIEXT}(\underline T)$ de biextensiones y, de hecho, un
\underline{automorfismo involutivo}:
\[
 {}^{s}({}^{s}E)=E.
\]
Todo enunciado acerca de bitorsores $E,\ldots$ permite, por tanto,
deducir un enunciado simétrico aplicándolo a $\,{}^{s}E,\ldots$.
Podemos llamarlo el enunciado obtenido \emph{intercambiando izquierda y
derecha}. En general solo explicitaremos uno de los dos enunciados
simétricos, eligiendo entre izquierda y derecha, por ejemplo, según
nuestras preferencias políticas.
